\documentclass[11pt]{article}
% Shared local declarations. Document roots use the repository's business-admin family.
\usepackage{business-admin}
\usepackage{amsmath,amssymb,tabularx,longtable,array,booktabs,enumitem}
\usepackage{ragged2e}
\newcolumntype{Y}{>{\RaggedRight\arraybackslash}X}
\newcommand{\SuiteMeta}[2]{\def\DocTitle{#1}\def\DocSubtitle{#2}\def\DocVersion{1.0}\def\DocStatus{Educational and operational reference}\def\DocOwner{Parliamentary Procedure and Governance Suite}\def\DocSummary{Independently authored reference and working materials, based on RONR 12th edition.}}
\newcommand{\SuiteOpen}{\makedoctitle\vspace{-0.5em}\noindent\rule{\linewidth}{0.4pt}\par\medskip}
\newcommand{\RonrAnchor}[1]{\par\smallskip{\footnotesize\textbf{Primary reading:} Robert's Rules of Order Newly Revised, 12th ed., #1. Original synthesis; not official interpretation.}\par}
\newcommand{\WriteLine}{\par\noindent\rule{\linewidth}{0.35pt}\par\vspace{0.5cm}}
\newcommand{\WriteLines}[1]{\par\noindent\rule{\linewidth}{0.35pt}\par\vspace{0.35cm}\noindent\rule{\linewidth}{0.35pt}\par\vspace{0.35cm}\noindent\rule{\linewidth}{0.35pt}\par\vspace{0.25cm}}
\newcommand{\Field}[1]{\par\medskip\noindent\textbf{#1}\par\WriteLines{3}}
\newcommand{\GovCheckbox}{\(\square\)}
\setlength{\emergencystretch}{3em}

\SuiteMeta{Voting, Counting and Election Controls}{Technical reference | original application synthesis}
\begin{document}
\SuiteOpen
\section{Scope and interpretation}
This is an original operational reference based on RONR (12th ed.). It is subordinate to applicable law and the organization’s validly adopted governing documents. Rule citations identify the primary source; the authors of this suite have not issued official interpretations.

\section{Voting bases}
For an unqualified majority vote, more than half of votes cast by eligible voters is needed; abstentions and blank ballots are excluded from that denominator. Two-thirds requires at least two thirds of votes cast. A majority of members present and a majority of entire membership use different denominators.

\textbf{Practical audit hook.} Source anchor: RONR 44:1--10.

\section{Worked calculation}
Assume 21 eligible participants are present and 20 cast a yes/no vote: 11 yes and 9 no. A majority of votes cast passes (11>10), but a two-thirds vote fails (needs 14 of 20). A majority of members present passes (11>10.5); do not reuse the first denominator without checking the rule.

\textbf{Practical audit hook.} Evidence: total membership, attendance, eligible ballots, yes/no/abstentions/invalid, required threshold and announcement.

\section{Method and integrity}
Choose voice, rising, counted standing, show of hands where appropriate, ballot, or roll-call according to governing requirements. Preserve secrecy where prescribed and protect the member’s right to abstain. If an election produces no required majority, conduct additional balloting instead of declaring a plurality winner without an authorizing rule.

\textbf{Practical audit hook.} Source anchor: RONR 45--46.

\section{Tellers report}
Document candidates or options, counts, invalid/blank ballots according to applicable rules, necessary majority and chair announcement. Reconcile the number of ballots against eligible voters and avoid exposing secret individual choices.

\textbf{Practical audit hook.} Exercise: investigate one unexplained excess ballot without inventing a tally.

\section{Independent review checklist}
\begin{enumerate}[leftmargin=*,itemsep=4pt]
\item Find the controlling adopted text and note its revision date.
\item Specify the decision, eligible decision maker, notice, quorum and vote basis.
\item Record evidence without fabricating approvals or preserving private member data in a public repository.
\item Classify unresolved issues as parliamentary, legal, governance policy or factual disputes.
\end{enumerate}
\end{document}
